\section{Output scale and membership in the restricted prior}
\label{sec:output-rescaling}

The restriction is a property of a parameterized predictor. It is not a
property of its Boolean sign function alone. This matters when comparing
networks trained with different losses or readout conventions.

\begin{proposition}[Every strictly signed finite-domain predictor can be rescaled]
\label{prop:output-rescaling}
Fix a finite input domain and a tanh network whose clean outputs obey
$a:=\min_x|z_w(x)|>0$. Write
\[
 B_s=\prod_{j=s+1}^{H}\|W^j\|_{\rm op},\qquad
 B=c\max\left\{1,
                 \sqrt{\frac1H\sum_{s=1}^{H}B_s^2}\right\},
\]
where an empty product is one. Keep every hidden weight fixed and replace
the readout by $\lambda v$. If $\lambda a>1/2+B$, the resulting network
satisfies the global clipped-score condition with \emph{zero} score
change, at the radius prescribed for its new readout. It also satisfies
the weaker margin-band condition.
\end{proposition}

\begin{proof}
The $1$-Lipschitz property of tanh and layerwise propagation give, for
every simultaneous disturbance,
\[
 |z_{\lambda v}(x;\eta)-\lambda z_v(x)|
 \le \|\eta\|_*
       \sqrt{1+\lambda^2 E_v\sum_{s=1}^{H}B_s^2}.
\]
At the calibrated radius this is at most
\[
 c\sqrt{\frac{1+\lambda^2 E_v\sum_s B_s^2}
                 {1+\lambda^2 H E_v}}
 \le B.
\]
Thus clean and noisy outputs remain beyond the same clipping endpoint
at every input. Both bounded scores are constant there.
\end{proof}

This observation does not supply a polynomial sample theorem for arbitrary
small analog circuits. The sufficient $\lambda$ can be exponentially large
in depth, or large because of a tiny minimum output on the finite domain.
Its readout energy becomes $\lambda^2E_v$, and the Gaussian probability
cost is charged accordingly. The main theorem instead constructs a robust
target implementation with polynomial probability cost. The empirical
protocol retains the actual trained readout and does not rescale a saved
model to make the condition pass.

Likewise, a larger allowed inverse-polynomial factor in the noise radius
is a meaningful calibration choice only when it is stated. The theorem
allows such choices because it tracks polynomial dependence on $c^{-1}$;
it does not turn an exponentially small required noise radius into a
polynomial-depth guarantee.
